在真空区域 $z>0$ 中放置一点电荷 $q>0$，其位置为 $(0,0,h)$，$h>0$。无限延伸理想导体占据 $z\le0$，保持接地电势 $\phi=0$，导体内静电场为零。设柱坐标的径向距离为 $\rho=\sqrt{x^2+y^2}$，忽略电荷的运动及辐射。
(1) 使用镜像电荷构造 $z>0$ 的完整静电势，证明符合边界条件与 Poisson 方程；(2) 从界面法向电场求诱导面电荷密度 $\sigma(\rho)$ 并积分其总电荷；(3) 求导体对实际点电荷的垂直静电力，说明相互作用势能必须含镜像法的 $1/2$ 因子；(4) 计算将实际点电荷从有限高度 $h$ 准静态移至无穷远所需的外力功，并讨论 $h\to\infty$ 与 $h\to0^+$ 的适用极限。
